\renewcommand{\Anno}{2011}

\renewcommand{\NumeroProblema}{3}

\renewcommand{\Testo}{ 
Paolo dispone di un lungo bastone cilindrico a sezione circolare.
Sezionando il bastone mediante opportuni piani, Paolo ricava tre pezzi
che poi, ricomposti opportunamente, costituiscono una cornice a forma
di triangolo equilatero.

Indichiamo con $r$ il raggio della sezione del bastone, e con $L$ il 
lato del triangolo equilatero i cui lati stanno sugli assi dei 3 
pezzi di bastone che costituiscono la cornice.

Determinare, in funzione di $r$ ed $L$, il volume e l'area della 
superficie della cornice.
}

\renewcommand{\Risposta}{
$\mbox{Volume}=3\pi Lr^{2}$, $\mbox{Area superficie}=6\pi Lr$.
}

\renewcommand{\Aiutino}{
Smontando i 3 ``lati'' della cornice, e rimontandoli opportunamente, 
si ottiene un cilindro~\ldots
}

\renewcommand{\SoluzioneUfficiale}{\TitoloSoluzione
Il volume $V$ e l'area $S$ della superficie della cornice sono uguali 
a
$$V=3\pi Lr^{2},
\hspace{4em}
S=6\pi Lr.$$

L'idea fondamentale \`{e} che, smontando i 3 pezzi che compongono la 
cornice, e rimontandoli opportunamente, si ottiene un cilindro di 
altezza $3L$ e sezione di raggio $r$. A quel punto il volume e la 
superficie della cornice coincidono, rispettivamente, con il volume e 
la superficie laterale del cilindro, da cui le formule di sopra.

La scomposizione ed il successivo rimontaggio sono illustrati nella 
figura seguente.

\begin{center}
	\psset{unit=2.5ex}
	\pspicture(-8,-5.5)(8,8)
	\pstGeonode[PointName=none,PointSymbol=none](4;90){A}(4;210){B}(4;330){C}
	\pstGeonode[PointName=none,PointSymbol=none](5.5;90){A'}(5.5;210){B'}(5.5;330){C'}
	\pstGeonode[PointName=none,PointSymbol=none](7;90){A''}(7;210){B''}(7;330){C''}
	\pspolygon[linewidth=1.5\pslinewidth](A)(B)(C)
	\pspolygon[linecolor=cyan](A')(B')(C')
	\pcline[linecolor=cyan]{<->}(B')(C')
	\lput*{:U}{$L$}
	\pspolygon[linewidth=1.5\pslinewidth](A'')(B'')(C'')
	\psline(A)(A'')
	\psline(B)(B'')
	\psline(C)(C'')
	\endpspicture
	
	\pspicture(-3,-3)(29,1)
	\pspolygon*[linecolor=green](0,0)(0,-1.5)(!-1.5 3 sqrt mul -1.5)
	\pspolygon(0,0)(!15 3 sqrt mul 0)(!16.5 3 sqrt mul -1.5)(!-1.5 3 sqrt 
	mul -1.5)
	\psline(!4 3 sqrt mul 0)(!5.5 3 sqrt mul -1.5)
	\psline(!11 3 sqrt mul 0)(!9.5 3 sqrt mul -1.5)
	\psline[linecolor=cyan](!-0.75 3 sqrt mul -0.75)(!15.75 3 sqrt mul -0.75)
	\psline[linestyle=dashed](0,1)(0,-2.5)
	\endpspicture

	\pspicture(-3,-3)(29,1)
	\pspolygon*[linecolor=green](!15 3 sqrt mul 0)(!16.5 3 sqrt mul 0)
	(!16.5 3 sqrt mul -1.5)
	\pspolygon(0,0)(!16.5 3 sqrt mul 0)(!16.5 3 sqrt mul -1.5)(0,-1.5)
	\psline(!4 3 sqrt mul 0)(!5.5 3 sqrt mul -1.5)
	\psline(!11 3 sqrt mul 0)(!9.5 3 sqrt mul -1.5)
	\psline[linecolor=cyan](0,-0.75)(!16.5 3 sqrt mul -0.75)
	\psline(!15 3 sqrt mul 0)(!16.5 3 sqrt mul -1.5)
	\pcline[linecolor=cyan,offset=-2ex]{|<->|}(0,-1.5)(!16.5 3 sqrt mul -1.5)
	\lput*{:U}{$3L$}
	\pcline[linecolor=cyan,offset=-2ex]{|<->|}(0,0)(0,-0.75)
	\Bput{$r$}
	% \psline[linestyle=dashed](!16.5 3 sqrt mul 1)(!16.5 3 sqrt mul -2.5)
	\endpspicture
\end{center}

Nell'ultimo passaggio abbiamo ``tagliato un angolo'' per ``rimontarlo 
dall'altra parte''. Si noti che i 3 pezzi della cornice sono ottenuti
tagliando mediante piani che formano angoli di $30^{\circ}$ con 
l'asse del cilindro, ma questo non \`{e} rilevante ai fini del 
successivo rimontaggio.

}

\renewcommand{\SoluzioniAlternative}{
}

\renewcommand{\Osservazioni}{
}

\renewcommand{\Commenti}{
}
